\documentclass{article}

\usepackage{fontspec}
\setsansfont{CMU Sans Serif}%{Arial}
\setmainfont{CMU Serif}%{Times New Roman}
\setmonofont{CMU Typewriter Text}%{Consolas}
\defaultfontfeatures{Ligatures={TeX}}
\usepackage[math-style=TeX]{unicode-math}
\usepackage[english, russian, ukrainian]{babel}

\usepackage[%
	a4paper,%
	footskip=1cm,%
	headsep=0.3cm,%
	top=2cm, %поле сверху
	bottom=2cm, %поле снизу
	left=2cm, %поле ліворуч
	right=2cm, %поле праворуч
    ]{geometry}


\usepackage{titlesec}

% Настройка стиля разделов
\titleformat{\section}
    {\normalfont\normalsize\bfseries} % стиль заголовка (здесь полужирный)
    {Завдання \thesection.}                    % формат номера (с точкой)
    {0.5em}                           % расстояние между номером и текстом
    {}                                 % код перед заголовком

\title{Питання та задачі з фізики плазми.}
\author{}
\date{}

\begin{document}

\maketitle

\section{Теоретичні питання}

\begin{enumerate}
    \item Які сили з боку магнітного поля можуть діяти на плазму?
    \item Які фізичні причини обумовлюють можливість утримання плазми магнітним полем?
    \item Які конфігурації магнітного поля здатні забезпечити утримання плазми?
    \item Які конфігурації магнітного поля здатні забезпечити утримання аксіально-симетричного плазмового стовпа?
    \item В чому полягає явище діамагнетизму плазми? Як воно впливає на утримання плазми магнітним полем?
    \item Які сили утримують плазмовий стовп у $Z$-пінчі?
    \item В чому полягає фізичний зміст умови Беннета?
    \item Чому в умову Беннета не входить радіус плазмового стовпа?
    \item Чому тороїдальної компоненти магнітного поля недостатньо для утримання тороїдального плазмового шнура?
    \item Що являють собою альвенівські хвилі? Від чого залежить їхня фазова швидкість?
    \item Опишіть особливості магнітозвукових хвиль у плазмі з малим $\beta$ порівняно з іншими типами низькочастотних хвиль анізотропної плазми.
    \item Що являють собою швидка та повільна магнітозвукові хвилі?
    \item Намалюйте картину магнітних силових ліній у швидкій магнітозвуковій хвилі, що поширюється перпендикулярно до магнітного поля.
\end{enumerate}

\section{Задачі}

\begin{enumerate}
    \item Яка індукція магнітного поля потрібна, щоб урівноважити тиск плазми в $1$ Торр?
    \item Який струм повинен протікати вздовж плазмового стовпа радіусом 10 см, щоб створюване цим струмом магнітне поле його утримувало? Концентрація плазми -- $10^{10}$ см$^{-3}$, температура -- $10$ еВ.
    \item Магнітне поле індукцією $10$ кГс, кінці силових ліній якого вморожені в провідні стінки, розташовані на віддалі 30 см одна від одної, утримує нижню межу повністю іонізованої цезієвої плазми з концентрацією $10^{12}$ см$^{-3}$. Оцініть розмір перехідної області між плазмою та магнітним полем.
\end{enumerate}

\end{document}